% chapters/ch05-normalization.tex
\chapter{Redundancy and Normalization}

\begin{learningobjectives}
In this chapter you will learn why redundancy is dangerous and how normalization helps.
You will be able to recognise update, insert, and delete anomalies, explain the idea of functional dependency in plain language, and apply a minimal normalization toolkit (1NF, 2NF, 3NF) without turning the topic into algebra.
\end{learningobjectives}

\section{Why redundancy is a correctness problem}
Redundancy looks harmless at first.
You store the same fact in multiple places, and everything still works.
The problem is that redundant facts do not remain consistent by accident.
They remain consistent only if every future update is done perfectly.
Real systems do not behave that way, especially when multiple people and services write data.

Normalization is not primarily about aesthetics.
It is a way of designing schemas so that each fact has a clear home.
When facts have a clear home, you reduce inconsistency and simplify enforcement of rules.

\begin{definitionbox}
\term{Redundancy} is the repeated storage of the same fact in multiple places.
A redundancy becomes harmful when it creates \term{anomalies}: situations where valid updates, inserts, or deletes produce invalid states or unintended data loss.
\end{definitionbox}

\begin{principlebox}
\textbf{One fact, one home}\\
If two rows can disagree about the same fact, the schema is already failing.
Normalization is the discipline of giving each fact a single home.
\end{principlebox}

\section{How to spot redundancy early}
Redundancy shows up in predictable ways.
If you learn to recognise them early, you avoid large refactors later.

Repeated text fields are the first signal.
If the same name, address, or description appears in many rows, ask which entity actually owns that fact.

Code+name duplication is another signal.
If you store a code and its label together across many rows, you are likely repeating a fact that should live in its own relation.

Copy-paste facts are the most dangerous signal.
If you find yourself copying attributes from one table to another “for convenience,” you are creating a future inconsistency.

\begin{figure}[htbp]
\centering
\includegraphics[width=0.85\textwidth]{normalization-before-after.pdf}
\caption{Normalization separates entity facts from event facts to remove redundancy.}
\end{figure}

\section{The three classic anomalies}
The anomaly triad is worth learning because it gives you a practical detection method.

An \term{update anomaly} happens when a fact appears in several rows and must be updated in all of them.
If one row is missed, the database ends up containing contradictory facts.

An \term{insert anomaly} happens when you cannot insert one fact without inventing another.
This often appears when a schema forces you to create a relationship record just to represent an entity that should exist independently.

A \term{delete anomaly} happens when deleting one record removes a fact you did not intend to remove.
The schema couples facts that should be independent.

These anomalies are not rare edge cases.
They happen naturally in schemas that mix “entity facts” with “event facts” in the same structure.

\section{Normalization reasoning and worked examples}
Normalization is easiest to understand when you start from something that looks like a spreadsheet.
We will use informal dependencies, a minimal toolkit, and two worked examples to show the reasoning.
Consider a naive attempt to record appointments:

\begin{examplebox}
\textbf{Bad design: mixed facts in one relation}\par\smallskip
\[
\begin{aligned}
\rel{AppointmentRecord}(&\underline{appointment\_id}, start\_time, doctor\_name, \\
&doctor\_specialty, patient\_name, patient\_birth\_date)
\end{aligned}
\]
\end{examplebox}

This table mixes appointment facts (start\_time) with doctor facts (name, specialty) and patient facts (name, birth\_date).
As a result, the same doctor information will be repeated in every row for that doctor.
The same patient information will be repeated in every row for that patient.

Now the anomalies appear immediately.

If a doctor changes specialty, you must update every appointment row for that doctor.
That is an update anomaly.
If you want to register a doctor before any appointments exist, you cannot do it without creating a fake appointment record.
That is an insert anomaly.
If you delete the only appointment a patient has, you may lose the only stored copy of that patient’s details.
That is a delete anomaly.

The table is not merely “not normalized”.
It makes correct operations fragile.

\subsection{The informal idea of functional dependency}
Normalization is often taught with formal definitions.
You can do that later.
For design work, you mainly need one idea: some attributes determine others.

If knowing one attribute always implies a specific value of another attribute, then the first attribute determines the second.
For example, if patient\_id uniquely identifies a patient, then patient\_id determines patient\_birth\_date.
In other words, patient facts should live where patient\_id is the key.

\begin{definitionbox}
An \term{informal functional dependency} means: “If you know \(X\), you can determine \(Y\).”
We write this as \(X \rightarrow Y\).
\end{definitionbox}

In the bad appointment table, we implicitly claim:
\[
appointment\_id \rightarrow doctor\_name,\; doctor\_specialty,\; patient\_name,\; patient\_birth\_date
\]
But that is conceptually wrong.
An appointment id should determine appointment details.
Doctor details should be determined by doctor identity, and patient details by patient identity.

Normalization is the act of aligning “what determines what” with keys and relation boundaries.

\begin{examplebox}
\textbf{University: course and department facts}\par\smallskip
If each course belongs to one department, then department data should not be repeated in the course offering rows:
\[
\begin{aligned}
\rel{Course}(&\underline{course\_id}, code \;[\text{unique}], title, \\
&department\_id \rightarrow \rel{Department})
\end{aligned}
\]
\[
\rel{Department}(\underline{department\_id}, name, office\_location)
\]
Course offerings can then reference \rel{Course} without repeating department facts.
\end{examplebox}

\subsection{A minimal normalization toolkit}
You do not need to memorise every normal form in the beginning.
A minimal toolkit is enough to produce robust schemas in most student projects.

\begin{principlebox}
\textbf{Normalize by dependency}\\
If an attribute depends on only part of a key, or on another non-key attribute, it belongs elsewhere.
Use dependencies to guide where facts live.
\end{principlebox}

\subsubsection{First Normal Form (1NF)}
1NF is about representation.
Attributes should be atomic, and relations should not contain repeating groups.
If you store multiple values in one field, the relational model cannot constrain or query them reliably.

A common violation is a column like \texttt{phones = "912..., 934..., 965..."}.
This is not a single attribute.
It is a collection hidden inside a string.

The correction is not “split the string later”.
The correction is to represent the collection as a relation.

\begin{examplebox}
\textbf{From multivalued attribute to relation}\\[2pt]
Instead of:
\[
\rel{Member}(\underline{member\_id}, name, phones)
\]
use:
\[
\rel{Member}(\underline{member\_id}, name)
\]
\[
\rel{MemberPhone}(\underline{member\_id} \rightarrow \rel{Member}, \underline{phone\_number})
\]
\end{examplebox}

\subsubsection{Second Normal Form (2NF)}
2NF matters when a relation has a composite primary key.
A relation violates 2NF when a non-key attribute depends on only part of that composite key.

This often happens in association relations where someone incorrectly stores attributes of one side.
For example, consider an association \rel{DoctorSpecialty}(doctor\_id, specialty\_id).
If you also store \texttt{doctor\_name} in this relation, then doctor\_name depends only on doctor\_id, not on the pair.
That creates redundancy and anomalies.

The correction is simple: move doctor attributes to \rel{Doctor}.

\subsubsection{Third Normal Form (3NF)}
3NF is about avoiding non-key attributes determining other non-key attributes.
A common pattern is storing a code and its description together in the same relation when the description is determined by the code.

For example, storing \texttt{specialty\_name} inside \rel{Doctor} when \texttt{specialty\_id} exists and determines the name.
That produces duplication and inconsistent spellings.

The correction is to create \rel{Specialty} and reference it.

\subsection{Clinic: normalized direction}
A normalized direction is to separate doctor and patient facts from appointment facts.

\begin{examplebox}
\textbf{Normalized direction (conceptual)}\par\smallskip
\[
\rel{Doctor}(\underline{doctor\_id}, name, specialty)
\]
\[
\rel{Patient}(\underline{patient\_id}, name, birth\_date)
\]
\[
\begin{aligned}
\rel{Appointment}(&\underline{appointment\_id}, start\_time, \\
&doctor\_id \rightarrow \rel{Doctor}, patient\_id \rightarrow \rel{Patient})
\end{aligned}
\]
\end{examplebox}

Now each fact has a clear home.
Doctor facts live in \rel{Doctor}, determined by \texttt{doctor\_id}.
Patient facts live in \rel{Patient}, determined by \texttt{patient\_id}.
Appointment facts live in \rel{Appointment}, determined by \texttt{appointment\_id}.
Redundancy is reduced, and the anomalies largely disappear.

\subsection{E-commerce: order lines mixing product facts}
Order lines often mix product facts that should live elsewhere.
Consider a naive \rel{OrderLine} that stores product name and category in every row.

\begin{examplebox}
\textbf{Order line with duplicated product facts}\par\smallskip
\[
\begin{aligned}
\rel{OrderLine}(&\underline{order\_id} \rightarrow \rel{Order}, \\
&\underline{product\_id}, product\_name, category, quantity, unit\_price)
\end{aligned}
\]
\end{examplebox}

If product\_name changes, you must update every order line that references it.
The fix is to separate product facts into \rel{Product} and keep only the reference in \rel{OrderLine}.

\section{Normalization boundaries}
Normalization can be taken too far, especially if taught as a checklist.
In practice, schemas are sometimes denormalized deliberately for performance or simplicity, but that should be done with clear justification and with compensating controls.

Start by removing the anomalies you can see.
If a change removes a real inconsistency risk, it is worth doing.
If you cannot name the anomaly removed, the change is likely cosmetic.

In a first design pass, you should prefer clarity and correctness.
If you later decide to denormalize, you should do it consciously and document which anomalies you are reintroducing and how you will manage them.

\begin{notebox}
A safe rule for early designs: normalize until the anomalies you can clearly see are removed.
If you cannot explain what anomaly a transformation fixes, you may be refactoring for style rather than correctness.
\end{notebox}

\section{Common mistakes}
Normalization mistakes usually come from mixing levels of meaning.

Mistake: Treating event data as if it were entity data.\\
Fix: Move dates and statuses to the event relation (Loan, Appointment, OrderLine).

Mistake: Stopping after the first split without checking dependencies.\\
Fix: Ask what determines each attribute and move those that depend on other non-key attributes.

Mistake: Over-normalizing without a reason.\\
Fix: Name the anomaly you are removing; if you cannot, keep the structure.

Mistake: Duplicating code+name pairs across rows.\\
Fix: Create a separate lookup relation and reference it.

\section{Mini checklist}
Before you declare a schema “normalized enough,” check:
\begin{itemize}[leftmargin=*,nosep]
  \item Does each fact have one clear relation where it is stored?
  \item Are multivalued attributes represented as separate relations?
  \item Do non-key attributes depend only on the full key?
  \item Are codes and their descriptions separated when a dependency exists?
  \item Can you explain which anomaly each split removes?
  \item Are event facts (dates, statuses, quantities) stored on event relations?
  \item Are you avoiding copy-paste attributes between relations?
\end{itemize}

\begin{chaptersummary}
Normalization is best understood as anomaly prevention.
Redundancy creates update, insert, and delete anomalies, which eventually produce incorrect states.
Functional dependency gives you the intuition for where facts belong: attributes should depend on the key and only the key.
A minimal toolkit (1NF, 2NF, 3NF) is enough to produce robust schemas for most domains.
\end{chaptersummary}

\begin{exercisebox}
Answer in full sentences.
\begin{enumerate}[label=\arabic*.]
  \item Give one example of each anomaly (update, insert, delete) in a domain you know.
  \item Show one example of a multivalued attribute and rewrite it as a separate relation (1NF).
  \item Create a small association relation (N:M) and explain why storing an entity attribute inside it violates 2NF.
  \item Describe one transitive dependency you might see in a real system and how you would eliminate it (3NF).
  \item For a chosen domain, identify one place where you might denormalize later and what risks that would introduce.
  \item For a clinic or library, sketch a table that mixes entity and event facts and explain the anomalies it creates.
  \item For a university domain, list two attributes that should move to a separate relation and explain the dependency.
\end{enumerate}
\end{exercisebox}

In the next chapter we move from redundancy to explicit constraints.
Normalization reduces anomalies, but constraints make the remaining rules enforceable.
